\FigureLayoutDeclare{img/2011/zhejiang_liberal/q22_fig1.png}{5.5cm}{0bp 96bp 23bp 36bp}
\FigureLayoutDeclare{img/2011/zhejiang_liberal/q18_fig1.png}{4.0cm}{0bp 0bp 0bp 0bp}

\chapter{2011年浙江卷（文）}

\section{单选题}

\begin{problem}
若 \(P = \{x \mid x < 1\}\)，\(Q = \{x \mid x > -1\}\)，则（\quad）

\choices
    {\(P \subseteq Q\)}
    {\(Q \subseteq P\)}
    {\(\complement_{\R} P\subseteq Q\)}
    {\(Q \subseteq \complement_{\R} P\)}
\end{problem}

\begin{answer}
C
\end{answer}

\begin{solution}
由 \(P=\{x\in\R\mid x<1\}\)，得 \(\complement_{\R}P=\{x\in\R\mid x\geqslant 1\}\). 任取 \(x\in\complement_{\R}P\)，都有 \(x\geqslant 1>-1\)，所以 \(x\in Q\)，从而 \(\complement_{\R}P\subseteq Q\). 另外，取 \(x=-2\) 可知 \(P\not\subseteq Q\)，取 \(x=2\) 可知 \(Q\not\subseteq P\)，且 \(0\in Q\) 但 \(0\notin\complement_{\R}P\)，故只有第三个选项 正确.
\end{solution}

\begin{problem}
若复数 \(z = 1 + \i\)，\(\i\) 为虚数单位，则 \((1 + z)\cdot z =\)（\quad）

\choices
    {\(1 + 3\i\)}
    {\(3 + 3\i\)}
    {\(3 - \i\)}
    {\(3\)}
\end{problem}

\begin{answer}
A
\end{answer}

\begin{solution}
因为 \(1+z=2+\i\)，所以 \((1+z)z=(2+\i)(1+\i)=2+2\i+\i+\i^2=1+3\i\)，故第一个选项 正确.
\end{solution}

\begin{problem}
若实数 \(x\)，\(y\) 满足不等式组 \(\begin{cases}
x + 2y - 5 \geqslant 0,\\
2x + y - 7 \geqslant 0,\\
x \geqslant 0,\\
y \geqslant 0,
\end{cases}\) 则 \(3x + 4y\) 的最小值是（\quad）

\choices
    {\(13\)}
    {\(15\)}
    {\(20\)}
    {\(28\)}
\end{problem}

\begin{answer}
A
\end{answer}

\begin{solution}
可行域由 \(x+2y\geqslant 5\)，\(2x+y\geqslant 7\)，\(x\geqslant 0\)，\(y\geqslant 0\) 确定. 目标函数的等值线斜率为 \(-\frac34\)，沿可行域向左下方移动时，最小值出现在下边界的顶点. 两条斜边界的交点满足 \(x+2y=5\)，\(2x+y=7\)，解得 \((x,y)=(3,1)\). 可行域下边界的另外两个顶点为 \((5,0)\) 和 \((0,7)\). 于是 \(3x+4y\) 在这三个点的值分别为 \(13\)，\(15\)，\(28\)，最小值为 \(13\)，故第一个选项 正确.
\end{solution}

\begin{problem}
若直线 \(l\) 不平行于平面 \(\alpha\)，且 \(l \not\subset \alpha\)，则（\quad）

\choices
    {\(\alpha\) 内的所有直线与 \(l\) 异面}
    {\(\alpha\) 内不存在与 \(l\) 平行的直线}
    {\(\alpha\) 内存在唯一的直线与 \(l\) 平行}
    {\(\alpha\) 内的直线与 \(l\) 都相交}
\end{problem}

\begin{answer}
B
\end{answer}

\begin{solution}
直线 \(l\) 不平行于平面 \(\alpha\)，所以 \(l\) 的方向不可能平行于平面 \(\alpha\) 内的直线. 因此，平面 \(\alpha\) 内不存在与 \(l\) 平行的直线. 又因为 \(l\not\subset\alpha\)，这不改变上述方向关系，故第二个选项 正确.
\end{solution}

\begin{problem}
在 \(\triangle ABC\) 中，角 \(A\)，\(B\)，\(C\) 所对的边分别为 \(a\)，\(b\)，\(c\). 若 \(a \cos A = b \sin B\)，则 \(\sin A \cos A + \cos^2 B =\)（\quad）

\choices
    {\(-\frac{1}{2}\)}
    {\(\frac{1}{2}\)}
    {\(-1\)}
    {\(1\)}
\end{problem}

\begin{answer}
D
\end{answer}

\begin{solution}
由正弦定理，存在 \(k>0\)，使 \(a=k\sin A\)，\(b=k\sin B\). 代入条件 \(a\cos A=b\sin B\)，得 \(k\sin A\cos A=k\sin^2B\)，从而 \(\sin A\cos A=\sin^2B\). 因此 \(\sin A\cos A+\cos^2B=\sin^2B+\cos^2B=1\)，故第四个选项 正确.
\end{solution}

\begin{problem}
若 \(a\)，\(b\) 为实数，则“\(0 < ab < 1\)”是“\(b < \frac{1}{a}\)”的（\quad）

\choices
    {充分而不必要条件}
    {必要而不充分条件}
    {充分必要条件}
    {既不充分也不必要条件}
\end{problem}

\begin{answer}
D
\end{answer}

\begin{solution}
若 \(0<ab<1\)，则 \(a\neq 0\). 当 \(a>0\) 时，除以 \(a\) 得 \(0<b<\frac1a\)，因而 \(b<\frac1a\)；当 \(a<0\) 时，除以 \(a\) 时不等号反向，得 \(\frac1a<b<0\)，此时 \(b<\frac1a\) 反而不成立，所以该条件不是充分条件. 反过来，取 \(a=1\)，\(b=-1\)，有 \(b<\frac1a\)，但 \(ab=-1\notin(0,1)\)，所以它也不是必要条件. 故第四个选项 正确.
\end{solution}

\begin{problem}
若某几何体的三视图如图所示，则这个几何体的直观图可以是（\quad）

\bitmapfigure[max width=0.42\linewidth]{img/2011/zhejiang_liberal/q7_fig1.png}

\choices
    {\choicebitmap[width=0.15\paperwidth]{img/2011/zhejiang_liberal/q7_optA.png}}
    {\choicebitmap[width=0.15\paperwidth]{img/2011/zhejiang_liberal/q7_optB.png}}
    {\choicebitmap[width=0.15\paperwidth]{img/2011/zhejiang_liberal/q7_optC.png}}
    {\choicebitmap[width=0.15\paperwidth]{img/2011/zhejiang_liberal/q7_optD.png}}
\end{problem}

\begin{answer}
B
\end{answer}

\begin{solution}
由三视图可知，正视图和俯视图均为带有水平分界线的矩形，而侧视图的上边界含有一条斜线. 逐一对照四个直观图，只有第二个选项在正面和顶面投影中同时形成相应的矩形分界，并在侧面投影中形成题图所示的斜边. 第一个选项和第三个选项会在正视图或俯视图中产生三角形轮廓，第四个选项的侧视图没有该斜边，故第二个选项正确.
\end{solution}

\begin{problem}
从装有 3 个红球，2 个白球的袋中任取 3 个球，则所取的 3 个球中至少有 1 个白球的概率是（\quad）

\choices
    {\( \frac{1}{10} \)}
    {\( \frac{3}{10} \)}
    {\( \frac{3}{5} \)}
    {\( \frac{9}{10} \)}
\end{problem}

\begin{answer}
D
\end{answer}

\begin{solution}
从 \(5\) 个球中任取 \(3\) 个球的取法数为 \(\mathrm C_5^3\). 所取的 \(3\) 个球中至少有 \(1\) 个白球的对立事件是取到 \(3\) 个红球，其取法数为 \(\mathrm C_3^3\). 所以所求概率为 \(1-\frac{\mathrm C_3^3}{\mathrm C_5^3}=1-\frac1{10}=\frac9{10}\)，故第四个选项 正确.
\end{solution}

\begin{problem}
已知椭圆 \(C_{1}: \frac{x^{2}}{a^{2}} + \frac{y^{2}}{b^{2}} = 1\) （\(a>b>0\)） 与双曲线 \(C_2: x^{2} - \frac{y^{2}}{4} = 1\) 有公共的焦点， \(C_2\) 的一条渐近线与以 \(C_1\) 的长轴为直径的圆相交于 \(A\)，\(B\) 两点. 若 \(C_1\) 恰好将线段 \(AB\) 三等分，则（\quad）

\choices
    {\(a^2 = \frac{13}{2}\)}
    {\(a^2 = 13\)}
    {\(b^{2} = \frac{1}{2}\)}
    {\(b^{2} = 2\)}
\end{problem}

\begin{answer}
C
\end{answer}

\begin{solution}
双曲线 \(C_2\) 的实半轴平方为 \(1\)，虚半轴平方为 \(4\)，所以其焦半距平方为 \(1+4=5\). 椭圆 \(C_1\) 与它有公共焦点，故 \(a^2-b^2=5\). 取 \(C_2\) 的渐近线 \(y=2x\)，它与圆 \(x^2+y^2=a^2\) 的交点横坐标为 \(\pm\frac{a}{\sqrt5}\). 由于 \(C_1\) 将线段 \(AB\) 三等分，且图形关于原点中心对称，椭圆与该直线的两个交点横坐标为 \(\pm\frac{a}{3\sqrt5}\). 将 \(y=2x\) 代入椭圆方程，得 \(\frac{x^2}{a^2}+\frac{4x^2}{b^2}=1\). 代入 \(x^2=\frac{a^2}{45}\)，得到 \(\frac1{45}+\frac{4a^2}{45b^2}=1\)，即 \(a^2=11b^2\). 与 \(a^2-b^2=5\) 联立，得 \(10b^2=5\)，所以 \(b^2=\frac12\)，故第三个选项 正确.
\end{solution}

\begin{problem}
设函数 \( f(x) = ax^2 + bx + c \)，（\(a\)，\(b\)，\(c\in\R\)）. 若 \(x = -1\) 为函数 \( y = f(x)\e^x \) 的一个极值点，则下列图象不可能为 \( y = f(x) \) 图象的是（\quad）

\choices
    {\choicebitmap[width=0.15\paperwidth]{img/2011/zhejiang_liberal/q10_optA.png}}
    {\choicebitmap[width=0.15\paperwidth]{img/2011/zhejiang_liberal/q10_optB.png}}
    {\choicebitmap[width=0.15\paperwidth]{img/2011/zhejiang_liberal/q10_optC.png}}
    {\choicebitmap[width=0.15\paperwidth]{img/2011/zhejiang_liberal/q10_optD.png}}
\end{problem}

\begin{answer}
D
\end{answer}

\begin{solution}
令 \(g(x)=f(x)\e^x\). 因为 \(\e^x>0\)，有 \(g'(x)=\e^x\bigl(f'(x)+f(x)\bigr)\). 由 \(x=-1\) 是极值点，首先有 \(g'(-1)=0\)，于是 \(f'(-1)+f(-1)=(-2a+b)+(a-b+c)=c-a=0\)，即 \(c=a\). 因此所有可能的抛物线都满足 \(f(x)=ax^2+bx+a\). 若第四个选项 的图象成立，则其开口向上，且两个零点都在 \(-1\) 的左侧. 设两个零点为 \(x_1\)，\(x_2\)，由韦达定理 \(x_1x_2=\frac{a}{a}=1\)，但 \(x_1<-1\)，\(x_2<-1\) 会推出 \(x_1x_2>1\)，矛盾. 因此第四个选项 的图象不可能.
\end{solution}

\section{填空题}

\begin{problem}
设函数 \( f(x) = \frac{4}{1 - x} \). 若 \( f(a) = 2 \)，则实数 \( a = \)\fillinblank{}.
\end{problem}

\begin{answer}
\(-1\)
\end{answer}

\begin{solution}
由 \(f(a)=2\)，得 \(\frac4{1-a}=2\). 由于右端非零，分母不为零，整理得 \(4=2-2a\)，所以 \(a=-1\)，且此时 \(1-a=2\neq0\)，满足定义域要求.
\end{solution}

\begin{problem}
若直线 \(x - 2y + 5 = 0\) 与直线 \(2x + my - 6 = 0\) 互相垂直，则实数 \(m = \)\fillinblank{}.
\end{problem}

\begin{answer}
\(1\)
\end{answer}

\begin{solution}
第一条直线可化为 \(y=\frac12x+\frac52\)，斜率为 \(\frac12\). 当 \(m\neq0\) 时，第二条直线的斜率为 \(-\frac2m\). 两直线垂直，所以 \(\frac12\left(-\frac2m\right)=-1\)，解得 \(m=1\). 当 \(m=0\) 时第二条直线为竖直线，与第一条直线不垂直，故答案为 \(1\).
\end{solution}

\begin{problem}
某小学为了解学生数学课程的学习情况，在 \(3000\) 名学生中随机抽取 \(200\) 名，并统计这 \(200\) 名学生的某次数学考试成绩，得到了样本的频率分布直方图（如图）. 根据频率分布直方图，\(3000\) 名学生在该次数学考试中成绩小于 \(60\) 分的学生数是\fillinblank{}.
\texfigure[max width=0.42\linewidth]{img_repaint/2011/zhejiang_liberal/q13_fig1.tex}
\end{problem}

\begin{answer}
\(600\)
\end{answer}

\begin{solution}
由直方图，成绩在 \(30\) 至 \(40\)，\(40\) 至 \(50\)，\(50\) 至 \(60\) 分组的频率分别为 \(10\times0.002=0.02\)，\(10\times0.006=0.06\)，\(10\times0.012=0.12\). 所以成绩小于 \(60\) 分的样本频率为 \(0.02+0.06+0.12=0.20\). 据此估计 \(3000\) 名学生中成绩小于 \(60\) 分的人数为 \(3000\times0.20=600\).
\end{solution}

\begin{problem}
某程序框图如图所示，则该程序运行后输出的 \(k\) 的值是\fillinblank{}.
\texfigure[max width=0.70\linewidth]{img_repaint/2011/zhejiang_liberal/q14_fig1.tex}
\end{problem}

\begin{answer}
\(5\)
\end{answer}

\begin{solution}
程序开始时 \(k=2\)，执行一次 \(k\leftarrow k+1\) 后得到 \(k=3\)，此时 \(4^3=64<3^4=81\)，继续循环. 随后 \(k=4\) 时 \(4^4=4^4\)，仍不满足 \(4^k>k^4\). 再执行一次循环得到 \(k=5\)，此时 \(4^5=1024>5^4=625\)，程序输出 \(k=5\).
\end{solution}

\begin{problem}
若平面向量 \(\bs{\alpha}\)，\(\bs{\beta}\) 满足 \(|\bs{\alpha}| = 1\)，\(|\bs{\beta}| \leqslant 1\)，且以向量 \(\bs{\alpha}\)，\(\bs{\beta}\) 为邻边的平行四边形的面积为 \(\frac{1}{2}\)，则 \(\bs{\alpha}\) 和 \(\bs{\beta}\) 的夹角 \(\theta\) 的取值范围是\fillinblank{}.
\end{problem}

\begin{answer}
\(\left[\frac{\pi}{6},\frac{5\pi}{6}\right]\)
\end{answer}

\begin{solution}
平行四边形的面积满足 \(|\bs{\alpha}|\,|\bs{\beta}|\sin\theta=\frac12\). 由 \(|\bs{\alpha}|=1\) 及 \(|\bs{\beta}|\leqslant1\)，得 \(\sin\theta=\frac1{2|\bs{\beta}|}\geqslant\frac12\). 面积为正，故 \(0<\theta<\pi\). 因此 \(\sin\theta\geqslant\frac12\) 等价于 \(\frac\pi6\leqslant\theta\leqslant\frac{5\pi}6\). 反之，对该区间内任意的 \(\theta\)，取 \(|\bs{\beta}|=\frac1{2\sin\theta}\leqslant1\) 即可满足面积条件，所以取值范围正是上述闭区间.
\end{solution}

\begin{problem}
若实数 \(x\)，\(y\) 满足 \(x^{2} + y^{2} + xy = 1\)，则 \(x + y\) 的最大值是\fillinblank{}.
\end{problem}

\begin{answer}
\(\frac{2\sqrt{3}}{3}\)
\end{answer}

\begin{solution}
令 \(s=x+y\)，\(d=x-y\). 则 \(x=\frac{s+d}{2}\)，\(y=\frac{s-d}{2}\)，从而 \(x^2+y^2+xy=\frac{3s^2+d^2}{4}=1\). 因此 \(3s^2\leqslant4\)，所以 \(s\leqslant\frac2{\sqrt3}=\frac{2\sqrt3}{3}\). 当且仅当 \(d=0\)，即 \(x=y\)，等号成立；此时取 \(x=y=\frac1{\sqrt3}\) 即满足原条件. 因此 \(x+y\) 的最大值为 \(\frac{2\sqrt3}{3}\).
\end{solution}

\begin{problem}
若数列 \(\left\{n(n + 4)\left(\frac{2}{3}\right)^n\right\}\) 中的最大项是第 \(k\) 项，则 \(k = \)\fillinblank{}.
\end{problem}

\begin{answer}
\(4\)
\end{answer}

\begin{solution}
记 \(u_n=n(n+4)\left(\frac23\right)^n\). 因为 \(u_n>0\)，考察相邻两项之比，得 \(\frac{u_{n+1}}{u_n}=\frac{(n+1)(n+5)}{n(n+4)}\cdot\frac23\). 当 \(n\in\N^*\) 时，\(\frac{u_{n+1}}{u_n}>1\) 等价于 \(2(n+1)(n+5)>3n(n+4)\)，即 \(n^2<10\). 所以 \(u_1<u_2<u_3<u_4\)，而从 \(n=4\) 开始相邻两项之比小于 \(1\)，数列递减. 故最大项为第 \(4\) 项，即 \(k=4\).
\end{solution}

\section{解答题}

\begin{problem}
已知函数 \( f(x) = A \sin \left( \frac{\pi}{3} x + \varphi \right) \)，\( x \in \R \)，\(A > 0\)，\(0 < \varphi < \frac{\pi}{2}\). \( y = f(x) \) 的部分图象如图所示，\(P\)，\(Q\) 分别为该图象的最高点和最低点，点 \(P\) 的坐标为 \((1,A)\).

\begin{enumerate}
\item 求 \( f(x) \) 的最小正周期及 \( \varphi \) 的值；
\item 若点 \(R\) 的坐标为 \((1,0)\)，\(\angle PRQ = \frac{2\pi}{3}\)，求 \(A\) 的值.
\end{enumerate}

\bitmapfigure[max width=0.38\linewidth]{img/2011/zhejiang_liberal/q18_fig1.png}
\end{problem}

\begin{solution}
\begin{enumerate}
\item 函数的角频率为 \(\frac\pi3\)，所以最小正周期 \(T=\frac{2\pi}{\pi/3}=6\). 又因为 \(P=(1,A)\) 是最高点，所以 \(\sin\left(\frac\pi3+\varphi\right)=1\). 由 \(0<\varphi<\frac\pi2\)，有 \(\frac\pi3<\frac\pi3+\varphi<\frac{5\pi}6\)，故只能取 \(\frac\pi3+\varphi=\frac\pi2\)，从而 \(\varphi=\frac\pi6\).
\item 由 \(\varphi=\frac\pi6\)，函数最低点的横坐标满足 \(\frac\pi3x+\frac\pi6=\frac{3\pi}2\)，结合图中所示位置得 \(Q=(4,-A)\). 于是 \(\overrightarrow{RP}=(0,A)\)，\(\overrightarrow{RQ}=(3,-A)\). 由夹角公式，
\[
\cos\angle PRQ=\frac{\overrightarrow{RP}\cdot\overrightarrow{RQ}}{|\overrightarrow{RP}|\,|\overrightarrow{RQ}|}=-\frac{A}{\sqrt{9+A^2}}
\]
代入 \(\angle PRQ=\frac{2\pi}{3}\)，得 \(-\frac{A}{\sqrt{9+A^2}}=-\frac12\)，所以 \(4A^2=9+A^2\). 因为 \(A>0\)，故 \(A=\sqrt3\).
\end{enumerate}
\end{solution}

\begin{problem}
已知公差不为 0 的等差数列 \(\{a_{n}\}\) 的首项 \(a_{1}\) 为 \(a\) （\(a \in \R\)），且 \(\frac{1}{a_1}\)，\(\frac{1}{a_2}\)，\(\frac{1}{a_4}\) 成等比数列.
\begin{enumerate}
\item 求数列 \( \{a_{n}\} \) 的通项公式；
\item 对 \(n \in \N^*\)，试比较 \(\frac{1}{a_2} + \frac{1}{a_{2^2}} + \frac{1}{a_{2^3}} + \cdots + \frac{1}{a_{2^n}}\) 与 \(\frac{1}{a_1}\) 的大小.
\end{enumerate}
\end{problem}

\begin{solution}
\begin{enumerate}
\item 设公差为 \(d\)，则 \(a_1=a\)，\(a_2=a+d\)，\(a_4=a+3d\). 因为三个倒数成等比数列，所以 \(\left(\frac1{a_2}\right)^2=\frac1{a_1a_4}\)，即 \(a(a+3d)=(a+d)^2\). 整理得 \(d(a-d)=0\). 由 \(d\neq0\)，得 \(d=a\)，于是 \(a_n=a+(n-1)d=na\). 由于 \(\frac1{a_1}\) 有意义，所以 \(a\neq0\).
\item 由 \(a_{2^j}=2^j a\)，所求和为 \(S_n=\frac1a\left(\frac12+\frac1{2^2}+\cdots+\frac1{2^n}\right)=\frac1a\left(1-\frac1{2^n}\right)\). 因此 \(S_n-\frac1a=-\frac1{a2^n}\). 当 \(a>0\) 时，\(S_n<\frac1a\)；当 \(a<0\) 时，\(S_n>\frac1a\).
\end{enumerate}
\end{solution}

\begin{problem}
如图，在三棱锥 \(P - ABC\) 中，\(AB = AC\)，\(D\) 为 \(BC\) 的中点，\(PO \perp\) 平面 \(ABC\)，垂足 \(O\) 落在线段 \(AD\) 上.
\begin{enumerate}
\item 证明： \(AP \perp BC\)；
\item 已知 \(BC = 8\)，\(PO = 4\)，\(AO = 3\)，\(OD = 2\). 求二面角 \(B - AP - C\) 的大小.
\end{enumerate}
\bitmapfigure[max width=0.36\linewidth]{img/2011/zhejiang_liberal/q20_fig1.png}
\end{problem}

\begin{solution}
\begin{enumerate}
\item 在等腰三角形 \(ABC\) 中，\(D\) 是 \(BC\) 的中点，所以 \(AD\perp BC\). 又因为 \(PO\perp\) 平面 \(ABC\)，所以 \(PO\perp BC\). 直线 \(AD\) 与 \(PO\) 相交且都在平面 \(APO\) 内，故 \(BC\perp\) 平面 \(APO\)，从而 \(AP\perp BC\).
\item 由 \(AO=3\)，\(OD=2\) 且 \(O\) 在线段 \(AD\) 上，得 \(AD=5\). 在等腰三角形 \(ABC\) 中，\(BD=DC=4\)，所以 \(AB^2=AD^2+BD^2=25+16=41\). 为计算二面角，建立直角坐标系，取 \(O\) 为原点，令 \(OD\) 为 \(x\) 轴正向，令 \(DB\) 为 \(y\) 轴正向的反向，令 \(OP\) 为 \(z\) 轴正向，则可取 \(A=(-3,0,0)\)，\(D=(2,0,0)\)，\(B=(2,-4,0)\)，\(C=(2,4,0)\)，\(P=(0,0,4)\). 于是 \(\overrightarrow{AP}=(3,0,4)\)，\(\overrightarrow{AB}=(5,-4,0)\)，\(\overrightarrow{AC}=(5,4,0)\). 在平面 \(PAB\) 和 \(PAC\) 内分别取
\[
\bs u=5\overrightarrow{AB}-3\overrightarrow{AP}=(16,-20,-12),\qquad \bs v=5\overrightarrow{AC}-3\overrightarrow{AP}=(16,20,-12)
\]
有 \(\bs u\cdot\overrightarrow{AP}=\bs v\cdot\overrightarrow{AP}=0\)，所以 \(\bs u\)，\(\bs v\) 分别是两个面内垂直于棱 \(AP\) 的直线方向向量. 又 \(\bs u\cdot\bs v=16^2-20^2+(-12)^2=0\)，故两条截线互相垂直，二面角 \(B-AP-C\) 的大小为 \(90^\circ\).
\end{enumerate}
\end{solution}

\begin{problem}
设函数 \(f(x) = a^2 \ln x - x^2 + ax\)，\(a > 0\).
\begin{enumerate}
\item 求 \( f(x) \) 的单调区间；
\item 求所有实数 \(a\)，使 \(\e - 1 \leqslant f(x) \leqslant \e^{2}\) 对 \(x \in [1, \e]\) 恒成立. 注：\(\e\) 为自然对数的底数.
\end{enumerate}
\end{problem}

\begin{solution}
\begin{enumerate}
\item 函数定义域为 \(x>0\)，且 \(f'(x)=\frac{a^2}{x}-2x+a=\frac{(a+2x)(a-x)}{x}\). 因为 \(a>0\)，\(x>0\)，所以 \(a+2x>0\)，导数符号由 \(a-x\) 决定. 因此 \(f(x)\) 在 \((0,a)\) 上单调递增，在 \((a,+\infty)\) 上单调递减.
\item 由 \(x=1\) 时的不等式，得 \(f(1)=a-1\geqslant\e-1\)，所以 \(a\geqslant\e\). 于是区间 \([1,\e]\) 位于 \((0,a]\) 内，函数在该区间上单调递增，最大值为 \(f(\e)=a^2-\e^2+a\e\). 要使上界成立，必须有 \(a^2+a\e-\e^2\leqslant\e^2\)，即 \((a-\e)(a+2\e)\leqslant0\). 结合 \(a>0\)，得到 \(a\leqslant\e\). 与 \(a\geqslant\e\) 联立，得 \(a=\e\). 此时 \(f(1)=\e-1\)，\(f(\e)=\e^2\)，且 \(f\) 在 \([1,\e]\) 上递增，故 \(a=\e\) 确实满足条件.
\end{enumerate}
\end{solution}

\begin{problem}
如图，设 \(P\) 为抛物线 \(C_1: x^2 = y\) 上的动点，过点 \(P\) 做圆 \(C_2: x^2 + (y + 3)^2 = 1\) 的两条切线，交直线 \(l: y = -3\) 于 \(A\)，\(B\) 两点.
\begin{enumerate}
\item 求 \(C_{2}\) 的圆心 \(M\) 到抛物线 \(C_{1}\) 准线的距离.
\item 是否存在点 \(P\)，使线段 \(AB\) 被抛物线 \(C_1\) 在点 \(P\) 处的切线平分，若存在，求出点 \(P\) 的坐标；若不存在，请说明理由.
\end{enumerate}
\bitmapfigure[max width=0.38\linewidth]{img/2011/zhejiang_liberal/q22_fig1.png}
\end{problem}

\begin{solution}
\begin{enumerate}
\item 抛物线 \(x^2=y\) 可写为 \(x^2=4py\)，所以 \(4p=1\)，其准线为 \(y=-\frac14\). 圆 \(C_2\) 的圆心为 \(M=(0,-3)\)，故 \(M\) 到准线的距离为 \(\left|-3+\frac14\right|=\frac{11}{4}\).
\item 设 \(P=(t,t^2)\)，其中 \(t\in\R\)，并设直线 \(l\) 上任意一点为 \(X=(u,-3)\). 过 \(P\)，\(X\) 的直线可写为 \((t^2+3)(x-u)-(t-u)(y+3)=0\). 圆心 \(M=(0,-3)\) 到此直线的距离为
\[
\frac{(t^2+3)|u|}{\sqrt{(t^2+3)^2+(t-u)^2}}
\]
该直线为圆 \(C_2\) 的切线，当且仅当上述距离等于半径 \(1\)，即
\[
(t^2+3)^2u^2=(t^2+3)^2+(t-u)^2
\]
记 \(K=(t^2+3)^2\)，整理得 \((K-1)u^2+2tu-(K+t^2)=0\). 设两条切线与 \(l\) 的交点横坐标为 \(u_1\)，\(u_2\)，由韦达定理，\(\frac{u_1+u_2}{2}=-\frac{t}{K-1}\)，所以线段 \(AB\) 的中点横坐标为 \(-\frac{t}{K-1}\).

当 \(t=0\) 时，抛物线在 \(P\) 处的切线为 \(y=0\)，与 \(l\) 平行，不能平分线段 \(AB\). 当 \(t\neq0\) 时，抛物线在 \(P\) 处的切线为 \(y=2tx-t^2\)，它与 \(l\) 的交点横坐标为 \(\frac{t^2-3}{2t}\). 因此平分条件为
\[
\frac{t^2-3}{2t}=-\frac{t}{K-1}
\]
因 \(K-1=(t^2+2)(t^2+4)\)，整理得
\[
(t^2-3)(t^2+2)(t^2+4)+2t^2=0
\]
令 \(s=t^2\)，则 \(s\geqslant0\)，上式化为 \(s^3+3s^2-8s-24=0\)，即 \((s+3)(s^2-8)=0\). 所以 \(s=2\sqrt2\)，从而 \(t=\pm\sqrt[4]{8}\). 对应的 \(y\) 坐标为 \(t^2=2\sqrt2\)，故存在这样的点，且
\[
P=\left(\sqrt[4]{8},2\sqrt2\right)\quad\text{或}\quad P=\left(-\sqrt[4]{8},2\sqrt2\right)
\]
这两个 \(t\) 均不为 \(0\)，且对应点 \(P\) 到圆心 \(M\) 的距离满足 \(PM^2=t^2+(t^2+3)^2>1\)，所以两条切线确实存在，所得点均符合题意.
\end{enumerate}
\end{solution}
